% ---------------------------------------------------------------------------
% editing_layer.tex -- the editing layer: equations, checks, and Figure S5.
% \input from simulator_figures.tex.  Code: src/14_editing.py (generator),
% src/15_validate_editing.py (checks), src/16_fig_editing.py (figures).
% Numbers: ../results/validate_editing.json and ../results/figS5_numbers.json.
% Theory record: notes/sciphy_notes.md S4, in particular S4.9.
% ---------------------------------------------------------------------------

% ===========================================================================
\section{The editing layer: what it simulates}\label{sec:edit}
% ===========================================================================

The tree layer says who descended from whom, and when. The editing layer says what each cell's tapes
read. It takes a stored tree from the library (\S\ref{sec:S1}) and writes edits along every branch.

\subsection*{Notation}

\begin{center}\small
\begin{tabular}{@{}lp{0.52\textwidth}l@{}}
\toprule
symbol & meaning & units / range \\
\midrule
$N$ & slots (sites) on one tape & $6$ \\
$k$, $z$ & tapes per cell; $z=1,\dots,k$ indexes them & count \\
$i$ & indexes symbols of the alphabet & --- \\
$t$ & calendar time as a fraction of the experiment ($T=1$ in the library) & $[0,1]$ \\
$e$ & a branch, from its parent's division at $t_e^0$ to its own division or
      harvest at $t_e^1$ (written $e$, since $b$ is the birth rate) & --- \\
$\lamz(t)$ & editing rate of tape $z$ & edits / tape / experiment \\
$\LamT$ & edits an average tape would receive over the whole experiment, if it never
          filled & edits / tape; $5.5$ \\
$r_z$ & speed of tape $z$ relative to the average & dimensionless, mean $1$ \\
$\lambda_0(t)$ & shape of the rate over time & dimensionless, mean $1$ \\
$\Lambda_z(t)$ & $\Lambda$-time: edits tape $z$ would have received by $t$ & edits / tape \\
$\mu_{e,z}$ & length of branch $e$ in $\Lambda$-time & edits / tape \\
$d$, $c$ & slots filled when a tape enters a branch; $c=N-d$ slots still open & $0$--$N$ \\
$j$ & slot index, $1$ written first; also ``the $j$-th edit on a tape'' & $1$--$N$ \\
$\nu$ & the $\nu$-th new edit on the current branch; it lands in slot $d+\nu$ & $1$--$c$ \\
$\kappa_g,\omega_g,W_g$ & knots, levels and cumulative weights of a piecewise-constant
      $\lambda_0$; epoch $g=1,\dots,G$ & dimensionless \\
$p_a(t)$ & share of edits going to signal channel $a$ & probability \\
$\xi^{(a)},\xi^{L}$ & symbol mix within signal channel $a$, and within the lineage
      channel; each sums to $1$ & probabilities \\
$\xi_i(t)$ & probability that an edit at time $t$ writes symbol $i$ & probability \\
$q_a$, $q_L$ & $\sum_i(\xi^{(a)}_i)^2$, $\sum_i(\xi^{L}_i)^2$: chance two independent writes
      inside one channel coincide & probability \\
$D$, $D_A$ & slots filled on a tape; at a common ancestor $A$ of two cells & $0$--$N$ \\
\bottomrule
\end{tabular}
\end{center}

\subsection{The model}\label{sec:edit-model}

A tape is a row of $N=6$ empty slots. Edits arrive at random times and fill the slots \emph{left to
right}, each writing one symbol. A full tape records nothing more. At a division both daughters
inherit a copy of the parent's tapes, and afterwards each fills its remaining slots independently.
Cells sharing an ancestor therefore share the edits written before they split: that is how tapes
record lineage. If the symbol mix changes over time, the slot a symbol sits in also says roughly
\emph{when} it was written. A tape is a queue, not a bag.

Formally this is SciPhy's recorder (\texttt{sciphy\_notes.md} \S1a), with the rate and the symbol mix
allowed to vary in time:
\clearpage
\begin{enumerate}
  \item While a tape has an open slot ($d<N$) it is edited as a Poisson process of
        rate $\lamz(t)$: edits arrive independently of one another and of what is already written.
  \item An edit writes symbol $i$ with probability $\xi_i(t)$ into slot $d+1$, and the
        depth becomes $d+1$.
  \item At $d=N$ the tape is full and stops recording.
  \item At a division both daughters start from an exact copy of the parent's
        tapes: same symbols, same depth.
  \item Given the tree, different tapes evolve independently.
\end{enumerate}
Everything below rests on the property in item~1: \emph{how fast an open tape is edited depends only
on the time, not on what is written on it.}

\paragraph{Tapes start empty at the clone founder.} Editing before the founding cell would leave a
prefix shared by every cell in the clone: it would consume slots while carrying no information
\emph{within} the clone. The simulator exposes it as $\Lambda_{\rm pre}$, the edits per tape spent
before $t=0$, set to $0$ throughout here. For a design study this is a definition rather than an
assumption --- the start of recording \emph{is} the start of the simulated experiment. For Park's data
it remains to be checked against the protocol.

\paragraph{Deliberately absent from this layer:} dropout (the next layer), heritable silencing
(row~A9), the separate mechanism at site~6, artefact symbols, and any coupling between the tapes of
one cell beyond the tree they share. A state-dependent rate is not built; the sampler only carries a
hook for it.

\subsection{The rate, and \texorpdfstring{$\Lambda$}{Lambda}-time}\label{sec:lambdatime}

\paragraph{Why a second clock.} When the rate varies, drawing edit times directly in calendar time is
awkward. Measuring time in \emph{expected edits} instead removes the difficulty: on that clock
editing always runs at exactly one expected edit per unit, however the real rate moves. We simulate
on the simple clock and convert each edit back afterwards.

\paragraph{The rate splits into three factors.}
\begin{equation}
  \lamz(t) \;=\; \underbrace{r_z}_{\text{this tape's speed}}\;\times\;
                 \underbrace{\LamT}_{\text{experiment total}}\;\times\;
                 \underbrace{\lambda_0(t)}_{\text{shape in time}},
  \qquad \int_0^1\!\lambda_0(u)\,du=1,\qquad \frac1k\sum_z r_z=1 .
  \label{eq:rate}
\end{equation}
$r_z$ carries the fact that tapes at different genomic sites edit at different speeds --- on Park's
Mouse3 the mean depth per tape runs from $1.33$ to $5.54$ slots. $\lambda_0$ carries any change over
time, such as a pause or a burst. Since $\lambda_0$ averages $1$, a tape with $r_z=1$ receives $\LamT$
expected edits over the experiment whatever the shape, so changing the shape never smuggles in a
change of total. Only $\lambda_0$ is varied in Fig.~\ref{fig:S5d}; $r_z=1$ everywhere in Fig.~S5.

\paragraph{``Averages to 1'' and ``integrates to 1'' are the same condition here.} The mean of a
function over an interval is its integral divided by the interval's length, and time runs over
$[0,1]$, so the two coincide. Were $T$ ever not $1$, they would differ by a factor $T$ and the
\emph{integral} is the condition that matters.

\paragraph{$\lambda_0$ is a relative rate, and is not bounded by $1$.} A value above $1$ means
``faster than the experiment average''. What the constraint fixes is the area, not the height: each
level is weighted by how long it lasts. In the schedule of Fig.~\ref{fig:S5d}, epochs of width
$0.25,0.10,0.25,0.40$ carry levels
\begin{equation}
  \lambda_0 = 1.18,\; 0,\; 1.57,\; 0.78,
  \qquad
  1.18(0.25)+0(0.10)+1.57(0.25)+0.78(0.40) = 1 ,
  \label{eq:lam0norm}
\end{equation}
so the burst runs at $1.57$, well above $1$: fast but short, against a last phase that is slow but
long. Multiplying by $\LamT=5.5$ gives the rates the figure plots, $6.47,0,8.63,4.31$ edits per tape
per experiment.

\paragraph{What is set, and what follows.} $\LamT$ is an \emph{input}, fixed before a run: Park's
three mouse arms measure $5.39$--$5.69$ edits per tape, so $\LamT=5.5$ here. The simulator never
infers it. The division of labour across eq.~\eqref{eq:rate} is then: $\LamT$ sets \emph{how much}
editing happens in total, $\lambda_0$ sets \emph{when} it happens and by its normalisation cannot
change the total, and $r_z$ sets how one tape compares with the average. That is why the two schedules
in Fig.~\ref{fig:S5d} end at the same depth: different $\lambda_0$, the same $\LamT$.

\paragraph{$\Lambda$-time} is the running total of the rate:
\begin{equation}
  \Lambda_z(t)=\int_0^t\!\lamz(u)\,du = r_z\LamT\,W(t),
  \qquad W(t)=\int_0^t\!\lambda_0(u)\,du .
  \label{eq:Lambda}
\end{equation}
Read $\Lambda_z(t)$ as ``how many edits tape $z$ would have taken by time $t$ had it never filled
up''; it only increases. \emph{$W$ is the quantity that runs from $0$ to $1$, not $\lambda_0$:} it is
the fraction of the experiment's total editing spent by time $t$, with $W(0)=0$ and $W(1)=1$. The two
are easy to confuse --- $\lambda_0(t)$ is the rate at an instant, mean $1$ and with no upper bound,
while $W(t)$ is its running total, increasing from $0$ to $1$. Substituting $\tau=\Lambda_z(t)$ turns a Poisson process of rate $\lamz(t)$
in $t$ into one of rate exactly $1$ in $\tau$, because $d\Lambda_z=\lamz\,dt$ makes the expected
number of edits in a stretch equal to that stretch's length in $\tau$. A branch $e$ then has a length
in $\Lambda$-time,
\begin{equation}
  \mu_{e,z}=\Lambda_z(t_e^1)-\Lambda_z(t_e^0)
  \qquad\text{(expected edits per tape during $e$, before saturation is applied).}
  \label{eq:mu}
\end{equation}
Every source of time variation --- an epoch boundary, a change of cell state, the end of a branch ---
then becomes the same kind of object: a point where the conversion between the two clocks changes
slope. One sampler handles them all, and the time-varying likelihood
(\texttt{sciphy\_notes.md} \S S4.4) is assembled from the same pieces.

\subsection{A rate that changes over time}\label{sec:piecewise}

\paragraph{The construction.} Give $\lambda_0$ a constant level inside each of a few epochs. Then
$\Lambda$ is a straight line within each epoch and converting back to calendar time is exact.
Take knots $0=\kappa_0<\kappa_1<\dots<\kappa_G=1$ and levels $\omega_g\ge0$ on
$[\kappa_{g-1},\kappa_g)$, rescaled so that $\sum_g\omega_g(\kappa_g-\kappa_{g-1})=1$ as
eq.~\eqref{eq:rate} requires. With cumulative weights
$W_g=\sum_{g'\le g}\omega_{g'}(\kappa_{g'}-\kappa_{g'-1})$, so $W_0=0$ and $W_G=1$, for $t$ in epoch
$g$:
\begin{equation}
  W(t)=W_{g-1}+\omega_g\,(t-\kappa_{g-1}),
  \qquad
  W^{-1}(w)=\kappa_{g-1}+\frac{w-W_{g-1}}{\omega_g}
  \quad\text{for } W_{g-1}\le w<W_g ,
  \label{eq:Winv}
\end{equation}
and $\Lambda_z^{-1}(\tau)=W^{-1}\!\big(\tau/(r_z\LamT)\big)$. The inverse is two steps: find which
epoch's interval $[W_{g-1},W_g)$ contains $w$ by comparing against the stored knot values, then invert
that epoch's straight line. No numerical root-finding.
\emph{Check.} At $G=1$, $\omega_1=1$: $\Lambda_z^{-1}(\tau)=\tau/(r_z\LamT)$, so a constant rate makes
$\Lambda$-time a rescaled calendar time.

\paragraph{A pause needs no special case.} A zero-rate epoch ($\omega_g=0$, for instance a Dox-gated
pause) has $W_{g-1}=W_g$. It occupies an empty interval of $\Lambda$-time, so the lookup can never
land inside it and no edit is ever placed there. Fig.~\ref{fig:S5d} shows exactly that.

\paragraph{What the construction would mean for inference.} Should the inference use the same
construction --- a skyline of epoch rates under a smoothing prior --- each $\omega_g$ is estimated as
an \emph{average} over its epoch. Resolution is then set by where the knots are put, not by the data,
and a sharp change returns as a ramp. A validation whose simulated truth shared the inference's knots
would flatter it, so later tests should plant change points that fall between them. Finely stepped
$\lambda_0$ approximates any smooth curve, so the generator itself does not restrict which truths can
be planted.

\subsection{Writing edits onto one branch}\label{sec:sampler}

\paragraph{The idea.} A tape enters branch $e$ with $d$ slots filled, so at most $c=N-d$ further edits
can be recorded there. On the $\Lambda$-clock the waits between successive edits are independent and
average one unit each. So draw $c$ waits, add them up into arrival times, and keep those that land
before the branch ends. Each kept arrival is an edit: convert its time back, draw its symbol, write it
into the next open slot.

\paragraph{The algorithm} (\texttt{decorate} in \texttt{14\_editing.py}), for one tape $z$ on one
branch $e$, branches taken parents before children:
\begin{enumerate}
  \item \textbf{Start.} The tape arrives at depth $d$, inherited from the end of the parent branch;
        the stem starts empty. Set $c=N-d$; if $c=0$ nothing happens here.
  \item \textbf{Waits.} Draw $E_1,\dots,E_c\sim\Exp(1)$, independently, in $\Lambda$-time units.
  \item \textbf{Arrivals.} $S_\nu=E_1+\dots+E_\nu$ for $\nu=1,\dots,c$: the $\Lambda$-time of the
        $\nu$-th new edit, measured from the start of the branch. These increase by construction, so
        no sorting is needed.
  \item \textbf{Which happened.} Keep those $\nu$ with $S_\nu\le\mu_{e,z}$ (eq.~\ref{eq:mu}); let
        $n_{\rm eff}$ be how many.
  \item \textbf{Back to calendar time.}
        $t_\nu=\Lambda_z^{-1}\big(\Lambda_z(t_e^0)+S_\nu\big)$, by eq.~\eqref{eq:Winv}.
  \item \textbf{Symbols.} Draw each edit's symbol from $\xi(t_\nu)$ (\S\ref{sec:channels}) and write
        it into slot $d+\nu$.
  \item \textbf{Hand-off.} The branch ends at depth $d+n_{\rm eff}$; both daughters inherit it.
\end{enumerate}

\paragraph{Why each step is right.}
\emph{Step 2:} in a Poisson process of rate $1$ the gaps between arrivals are independent exponentials
of mean $1$ --- this is what ``memoryless'' means. The draws are in $\Lambda$-time units precisely
because the process has rate $1$ there (\S\ref{sec:lambdatime}). At a constant calendar rate $\lambda$
the waits would be $\Exp(\lambda)$, mean $1/\lambda$; stretching time by $\lambda$ makes them $\Exp(1)$.
\emph{Step 4, and why $c$ draws are enough:} the real tape stops at its $c$-th edit, and stopping never
moves the arrivals before it, so the edits that occur are exactly the first $\min(\text{arrivals},c)$.
Drawing more would be a device, not extra biology.
\emph{Step 7:} depth has to be carried across branches because saturation couples a branch to
everything above it --- a tape that filled early has nothing left to record later.

\paragraph{$\mu$ is an expectation, not a ceiling.} The number of $S_\nu$ landing below $\mu_{e,z}$ is
$\Poisson(\mu_{e,z})$, limited only by the $c$ open slots. Several short waits fit easily inside one
branch: at $\mu=0.5$ edits per tape, a branch carries no edit with probability $e^{-0.5}=0.61$, one
with probability $0.30$, and two or more with probability $0.09$.

\paragraph{A worked conversion.} With a constant rate, $r_z=1$ and $\LamT=5.5$, we have
$\Lambda_z(t)=5.5\,t$. A branch beginning at $t_e^0=0.4$ starts at $\Lambda_z(0.4)=2.2$, so an edit at
$S_\nu=0.7$ happened at $t_\nu=(2.2+0.7)/5.5=0.527$. Under a piecewise-constant rate the same step runs
through eq.~\eqref{eq:Winv}: fast epochs are stretched on the $\Lambda$-clock, so edits landing in them
map back into a narrower calendar window, and slow epochs the reverse.

\paragraph{A second, independent implementation.} Route~(i) draws the count first,
$n\sim\Poisson(\mu_{e,z})$, scatters $n$ points $\Uniform(0,\mu_{e,z})$, sorts them and keeps the first
$\min(n,c)$. It is exact by conditional uniformity: given how many arrivals a Poisson process produced
in an interval, their times are uniformly scattered over it (\texttt{sciphy\_notes.md} \S S4.2). It
samples the same law by a different route, which is what makes it a useful cross-check
(\S\ref{sec:editval}). The production route above never draws more than $c\le6$ numbers per tape per
branch and needs no sort.

\subsection{Which symbol gets written: channels}\label{sec:channels}

\paragraph{The idea.} Symbols fall into groups. \emph{Lineage} symbols are many and diverse; their job
is to tell lineages apart. \emph{Signal} symbols, one or a few per recorder, are the ones whose share
rises when a pathway is on. The signal changes the \emph{share} of edits going to its own channel, not
how often tapes are edited, because ENGRAM raises a pegRNA's abundance rather than the prime editor's
activity (\texttt{sciphy\_notes.md} \S0.2b).

\paragraph{The construction.} Every symbol belongs to exactly one channel. With $A$ signal channels,
\begin{equation}
  \xi_i(t)=
  \begin{cases}
    p_a(t)\,\xi^{(a)}_i, & i \text{ in signal channel } a,\\[3pt]
    \Big(1-\sum_a p_a(t)\Big)\,\xi^{L}_i, & i \text{ a lineage symbol,}
  \end{cases}
  \qquad \sum_i \xi_i(t)=1 .
  \label{eq:channels}
\end{equation}
These are two cases, not a sum: for a given symbol only one line applies. What they say together is
that the channels compete for the same edits. Raising $p_a$ from $0.02$ to $0.30$ multiplies every
lineage symbol's probability by the same factor, $0.98\to0.70$, without altering the lineage mix
$\xi^L$ itself. So ``$30\%$ of edits'' means precisely this: while the signal is on, each edit writes
the signal symbol with probability $0.30$ and some lineage symbol with probability $0.70$. It is a
probability per edit, not an amount per slot.

\paragraph{Sampling} goes in two steps: first the channel, with probabilities $p_1(t),\dots,p_A(t)$
and $1-\sum_a p_a(t)$; then the symbol from that channel's own mix. The within-channel mixes each sum
to $1$ because they are \emph{conditional} on the channel; multiplying the two steps reproduces
eq.~\eqref{eq:channels} exactly. The factored form is the one that is built, because it keeps the knob
the signal moves ($p_a$) separate from the designed, fixed mixes.

\paragraph{The budget couples recorders.} Every edit is one write, so $\sum_a p_a(t)\le1$ always, and
in practice $\lesssim0.3$, the shared-tape ceiling (\texttt{sciphy\_notes.md} \S I.7.5). A second
recorder does not receive capacity of its own: it either splits the same share or eats into the lineage
channel. This is the design-sweep statement of the $(p,j)$ trade-off, and it is why the partition is
built into the data structure before it is ever varied.

\paragraph{The lineage mix used here} is Park's measured one --- the three mouse arms' symbol
frequencies pooled by edit count, $113$ symbols with $q_L=0.01664$, an effective alphabet
$1/q_L=60$. Only frequency values are read; the sequences themselves never leave the cluster.

\subsection{What a signal channel costs the lineage record}\label{sec:collision}

Two unrelated cells look related at a slot when they happen to write the same symbol there. A frequent
signal symbol makes that far more likely, and that is the price the lineage record pays for carrying a
signal. Two independent writes at times $t_1,t_2$ coincide with probability
$\langle\xi(t_1),\xi(t_2)\rangle=\sum_i\xi_i(t_1)\xi_i(t_2)$. Under eq.~\eqref{eq:channels}, with the
within-channel mixes fixed and only the shares moving, the sum splits by channel because the channels
hold disjoint symbols:
\begin{equation}
  \langle\xi(t_1),\xi(t_2)\rangle
  =\sum_a p_a(t_1)p_a(t_2)\,q_a
  +\Big(1-\sum_a p_a(t_1)\Big)\Big(1-\sum_a p_a(t_2)\Big)q_L .
  \label{eq:collision}
\end{equation}
The signal channel contributes $p_a(t_1)p_a(t_2)q_a$, largest when \emph{both} writes fall while it is
induced. With a single signal symbol ($q_a=1$) at $p=0.3$ against Park's $q_L$, the collision
probability is $0.09+0.49\times0.01664=0.0982$, which is $5.9$ times $q_L$: the effective alphabet drops
from $60$ to $10$. Validation reproduced this on simulated tapes, $0.0976$ observed against $0.0982$
predicted (\S\ref{sec:editval}). It is a property of the model, not a measurement on ENGRAM data.

\subsection{What is stored}\label{sec:storage}

An edit happens on exactly one branch and is inherited by every cell below it, so edits are stored
\emph{per branch}, once. For each branch and tape the record holds how many edits were written there,
which slot the first went into, and their symbols and calendar times --- at most $N$ of each. A cell's
tape is read by walking root to cell and joining the branch records in order. That walk uses only the
stored record, never the sampler's own hand-off bookkeeping, so a hand-off error surfaces as a mismatch
instead of being copied silently. The walk also notes which branch wrote each slot: two cells whose slot
shares an origin share that edit, which is what \S\ref{sec:editval} needs to avoid counting the same
comparison many times. Storage is at most $6nk$ edits per tree, about $360{,}000$ at $n=1{,}976$,
$k=30$.

\subsection{The closed forms the layer is checked against}\label{sec:editcf}

\paragraph{A --- depth anywhere on a lineage.} Along one root-to-cell path a tape sees a single Poisson
process in $\Lambda$-time, stopped at $N$; branches merely cut that path into pieces. So at time $t$,
writing $\mu=\Lambda_z(t)$,
\begin{equation}
  \Prob(D=d)=e^{-\mu}\frac{\mu^{d}}{d!}\ \ (d<N),\qquad
  \Prob(D=N)=1-\sum_{d<N}e^{-\mu}\frac{\mu^{d}}{d!},\qquad
  \E[D]=\sum_{j=1}^{N}\Prob\big(\Poisson(\mu)\ge j\big).
  \label{eq:depth}
\end{equation}
Differentiating the mean gives the rate at which a lineage accumulates filled slots,
\begin{equation}
  \frac{d}{dt}\,\E[D(t)] \;=\; \lamz(t)\,\Prob\big(\Poisson(\Lambda_z(t))\le N-1\big):
  \label{eq:depthslope}
\end{equation}
the editing rate times the chance the tape is still open. At a constant rate with $\LamT=5.5$ this
falls from $5.5$ slots per unit time at $t=0$ to $2.9$ at harvest, where $47\%$ of tapes are full.
\emph{Its blind spot:} eq.~\eqref{eq:depth} contains no tree at all, only elapsed $\Lambda$. It
therefore cannot notice an edit attached to the wrong lineage, which is why checks D and E below are
needed.

\paragraph{B --- which slot holds what.} Slot~1 is written first and slot~6 last, but each over a wide
window of times, so this is the only closed form that sees edit \emph{order}. Given that a tape ends at
depth $d$ with $\mu=\Lambda_z(1)$ elapsed, the $\Lambda$-time $x$ of its $j$-th edit has density
$K_{j,d}$:
\begin{itemize}
  \item \textbf{not full, $d<N$:} exactly $d$ edits fell in $[0,\mu]$, so by conditional uniformity
        their times are the order statistics of $d$ independent $\Uniform(0,\mu)$ draws, giving
        $x/\mu\sim\BetaDist(j,\,d+1-j)$;
  \item \textbf{full, $d=N$:} the $N$-th edit landed at or before $\mu$, so
        $K_{j,N}(x)\propto \gamma_j(x)\,\Gamma_{N-j}(\mu-x)$ on $0<x<\mu$, where $\gamma_j$ is the
        density of the $j$-th arrival (Gamma, shape $j$, rate $1$) and $\Gamma_{N-j}$ is the
        probability that the remaining $N-j$ arrivals still fit into what is left. At $j=N$ the second
        factor is $1$ and this is a Gamma of shape $N$ truncated to $(0,\mu)$.
\end{itemize}
Writing $\tilde\xi_i(x)=\xi_i\big(\Lambda_z^{-1}(x)\big)$ for the mix as a function of $\Lambda$-time,
the expected share of symbol $i$ at slot $j$ among tapes of final depth $d$ is
\begin{equation}
  \E\big[\text{share of } i \text{ at slot } j \mid d\big]
  = \int_0^{\mu} K_{j,d}(x)\,\tilde\xi_i(x)\,dx .
  \label{eq:slotshare}
\end{equation}
Under a constant mix this collapses to ``every slot shows $\xi$'', which is why only a time-varying
mix gives the check teeth. The kernel conditions on $(\mu,d)$, so it must be applied within strata of
$r_z$ and final depth; pooling across tape speeds mixes kernels and makes it uninterpretable.

\paragraph{D --- the inherited prefix is exact.} Let two cells have last common ancestor $A$, whose
tape sat at depth $D_A$. Every edit on the path down to $A$ is shared by construction, and edits after
$A$ are independent between the two lineages. So the two tapes agree \emph{exactly} in slots
$1,\dots,D_A$. This is an assertion with zero tolerance, not a statistic.

\paragraph{E --- collisions past the divergence.} At slot $D_A+1$, if both lineages filled it, the two
writes are independent draws at their own times $t_1,t_2>t_A$, so they match with probability
$\langle\xi(t_1),\xi(t_2)\rangle$ of eq.~\eqref{eq:collision}. Because edit times are stored, the
expectation can be evaluated on the \emph{realised} pairs rather than assumed, which is what turns
eq.~\eqref{eq:collision} from a derivation into a checked result.

% ===========================================================================
\section{Validation: does the code produce that model?}\label{sec:editval}
% ===========================================================================

Section~\ref{sec:edit} defines a process; validation (\texttt{15\_validate\_editing.py}) asks whether
the code generates it. A failure here is a bug, not a finding. Every entry is a deviation from a
closed form of \S\ref{sec:editcf} measured in standard errors, in the sense of eq.~\eqref{eq:z}, with
one difference that matters: \textbf{the standard error is taken across trees}, never across cells or
pairs, because cells within one tree share ancestry and are not independent.

\paragraph{Three configurations}, all at $n=210$, $k=30$, $N=6$, $\LamT=5.5$, tapes empty at the
founder, 200 library trees spread over $\capt\in\{0.0005,0.002,0.02,0.1,0.25\}$ and
$\turn\in\{0.3,0.7\}$:
\begin{center}\small
\begin{tabular}{@{}lp{0.68\textwidth}@{}}
\toprule
name & what it sets \\
\midrule
\textsc{const} & constant rate; one signal channel at a fixed share $p=0.05$ \\
\textsc{tv}    & rate levels $1.2/0/1.6/0.8$ on knots $0,0.25,0.35,0.6,1$ --- including a
                 \emph{zero-rate window} --- and share $0.02/0.30/0.02$ on knots $0,0.4,0.7,1$,
                 deliberately not aligned with the rate's knots \\
\textsc{rhet}  & \textsc{const} with per-tape speeds $r_z$ lognormal (sd $0.5$), rescaled to mean $1$ \\
\bottomrule
\end{tabular}
\end{center}

\paragraph{The reading rule, fixed before the run.} No-effect value is $|z|=0$. \textbf{Pass} means
the worst $|z|$ over every tested cell is at most $3.0$ \emph{and} there are no violations of the hard
assertions. A cell is tested only when its expected positive and negative counts each reach 100, so
that its standard error is estimable; excluded cells are reported rather than hidden. With about 220
cells, one $|z|$ above $3$ is expected roughly $0.6$ times by chance, so a single excursion is re-run
at a fresh seed before it counts as a defect; two, or a systematic pattern of signs within one check,
would be a defect.

\begin{center}\small
\begin{tabular}{@{}l p{0.58\textwidth} r r@{}}
\toprule
check & what it tests & cells & worst $|z|$ \\
\midrule
A    & tip depth against eq.~\eqref{eq:depth} \emph{(tree-blind)}                 & 21 & 1.48 \\
A$'$ & internal-node depth at $\Lambda_z(t_{\rm node})$, in three time bins       & 63 & 1.81 \\
A$_r$& per-tape mean depth under spread $r_z$ \emph{(added)}                      & 5  & 1.61 \\
B    & signal share by slot given final depth, eq.~\eqref{eq:slotshare}
       \emph{(the only order check)}                                              & 42 & 2.03 \\
B$_L$& lineage mix pooled over slots, in ten equal-mass groups                    & 20 & 2.26 \\
C    & route (i) against route (ii), paired per tree                              & 38 & 2.67 \\
E    & collisions past the ancestor, random cell pairs                            & 15 & 1.58 \\
E$_u$& the same over unique comparisons \emph{(added)}                            & 15 & 1.65 \\
D    & hand-off slot equals the recomputed depth; depth $\le N$; edit times inside their branch;
       times ordered along each tape; valid symbol ids; prefix identical to the ancestor's depth
       & --- & $0$ of $1.9{\times}10^{8}$ \\
\bottomrule
\end{tabular}
\end{center}

\paragraph{Result: pass.} 219 cells tested, none excluded, worst $|z|=2.67$, none above $3$; zero
violations in $1.9\times10^{8}$ hard-assertion comparisons; the stored record reloads identically.

\paragraph{Two results worth more than the pass itself.}
\begin{itemize}
  \item \textbf{The order check had something to catch, and held.} Under \textsc{tv}, depth-6 tapes
        carry the signal symbol in $3.2/7.2/14.0/19.9/20.6/15.6\%$ of slots 1--6 --- the on-window
        rising and falling along the queue --- and eq.~\eqref{eq:slotshare} tracks it to about
        $0.001$. Passing under both a time-varying and a constant mix is what rules out the failure
        mode that would have mattered: symbols written into slots in the wrong order, which check A
        cannot see at all.
  \item \textbf{Eq.~\eqref{eq:collision} is verified.} When both writes fall in the on-window
        ($p=0.30$, one signal symbol) two lineages collide at $0.0976$ against $0.0982$ predicted;
        with one write in the window, $0.0176$ against $0.0174$; with neither, $0.0165$ against
        $0.0164$.
\end{itemize}

\paragraph{Two changes to the validation, both disclosed.} Check~E as first written compared randomly
chosen pairs of cells. Most such pairs compare \emph{the same two edits}, because a slot written just
below the common ancestor is shared by every cell beneath it. Per tree the statistic then rests on a
handful of independent comparisons of a rare event (a match has probability about $0.018$), and its
$t$-statistic is not calibrated: on a 20-tree trial run it reached $|z|=4.30$, and its
minimum-count rule overstated the information about elevenfold ($115{,}828$ expected matches against
$10{,}111$ unique comparisons). E$_u$ counts each (origin branch of edit 1, origin branch of edit 2,
tape) once. Both are unbiased, since the weighting never depends on the symbols, and the verdict's
scope was widened to include E$_u$ \emph{before} the 200-tree run; E as first written also passes at
200 trees. The general lesson --- \emph{comparisons between cells are redundant wherever the cells
share the edit} --- applies to any later statistic built on simulated tapes, but not to check~B, whose
closed form is stated per cell and would be biased by deduplication. A$_r$ was added so that the
per-tape speeds, a code path the $r_z\equiv1$ runs never touch, is exercised at all.

\paragraph{The yardstick is about $10\%$ generous.} A null built by splitting the trees into two random
halves and comparing gives a spread of $0.90$ rather than $1$ over the 219 cells, worst $2.25$. The
across-tree standard error is therefore slightly over-estimated, so every $|z|$ above is if anything
$\sim10\%$ small: the checks are marginally less sensitive than nominal, not more permissive.

\paragraph{Cost.} $0.05$~s per tree at $n=210$, $k=30$; $0.45$~s (production route) and $0.53$~s
(route~i) at $n=1{,}976$, plus $0.09$~s to assemble tapes; peak memory $0.11$~GB. Editing is not the
expensive layer --- the tree layer is.

\paragraph{What passing does not establish.} A and A$'$ are tree-blind and B, B$_L$ are marginals;
only D and E see the joint structure. And a simulator is also a model: passing means the code
generates SciPhy's recorder, not that the recorder is right for Park. Dropout, heritable silencing and
the site-6 mechanism are absent by design, and no Park-calibrated run has been made --- per-tape rates
cannot be read off Park while depth~0 is censored and dropout correlates with depth, so editing and
dropout will have to be fitted jointly.

% ===========================================================================
\section{Figure S5 --- the editing layer, illustrated}\label{sec:S5}
% ===========================================================================

\emph{Illustrations, not tests.} The formal record is \S\ref{sec:editval}; here each panel shows one
thing the layer does, with the closed form drawn as a line over the simulated points. Script
\texttt{16\_fig\_editing.py}; numbers in \texttt{../results/figS5\_numbers.json}. Unless a panel says
otherwise: 200 library trees of $n=210$ cells over the same $\capt$ and $\turn$ spread as
\S\ref{sec:editval}, $k=30$ tapes, $\LamT=5.5$, tapes empty at the founder, standard errors across
trees.

\subsection{What each panel shows}

\paragraph{(a) Inheritance.} One 8-cell tree with three tapes per cell drawn as rows of six slots,
every edit shaded by how many cells share it. Edits written on the stem fill the first slots of every
cell; later branches add lighter, private endings. This is the whole lineage-recording mechanism in
one picture, and it is the panel to show first.

\paragraph{(b) Relatedness is read off the shared prefix.} How many slots two cells share, against the
time their lineages split. The line is the \emph{full} expectation: the slots filled before the split,
eq.~\eqref{eq:depth}, plus the matches that happen by chance afterwards. Slot $D_A+m$ can only match by
chance if every earlier one did, so
\begin{equation}
  \E[\text{shared}] = \E[D_A] + \sum_{m\ge1} q^{\,m}\,
     \Prob\big(\text{both lineages write}\ge m\text{ more}\big)^2 ,
  \label{eq:shared}
\end{equation}
with $q=q_L$ here since no signal channel is switched on in (a) or (b). The chance term contributes at
most $0.017$ slots. Comparisons are taken one per (split, tape) --- one cell drawn from each side of
the split --- following the lesson of \S\ref{sec:editval}.

\paragraph{(c) Time is written into slot order.} A signal pulse, share $0.02\to0.30\to0.02$, on between
$t=0.4$ and $0.7$, at a \emph{constant} editing rate so that slot order reads directly as time order.
Read from tapes that filled all six slots. The top row is $K_{j,6}$ of eq.~\eqref{eq:slotshare}
expressed in calendar time --- when each slot was written --- and the bottom row is the resulting
signal share per slot against eq.~\eqref{eq:slotshare}.

\paragraph{(d) The rate can be programmed.} The rate schedule of \textsc{tv}, with its pause and burst.
The top row recovers the rate that the tapes actually experienced; the bottom row shows what that does
to the tape's depth. Note what the top row requires: dividing edits by the tape-time spent
\emph{open}. Dividing instead by all tape-time, open or full, produces a curve that sags although the
rate is flat, which is the distinction eq.~\eqref{eq:depthslope} makes between the rate and the tape
filling up.

\begin{figure}[p]
  \centering
  \includegraphics[width=\textwidth]{figS5a_inheritance}
  \caption{\textbf{Cells inherit their ancestors' edits as a shared prefix.} One library tree, $n=8$
  captured cells at $\capt=0.25$, $\turn=0.5$, with $k=3$ tapes drawn per cell. Horizontal axis is time
  since the founding cell. Dots are edits, placed at the time they were written and on the branch that
  wrote them; boxes to the right are each cell's three tapes at harvest, slot~1 leftmost. Shading is
  the same in both: how many of the eight cells carry that edit, from one cell (lightest) to all eight
  (darkest); white boxes are slots that were never filled. 33 edits were written on this tree.
  \emph{How to read it:} the dark edits on the stem appear in the same early slots of every cell,
  because every cell descends from the branch that wrote them. Reading a tape left to right is reading
  down the tree: shared ancestry first, private history last. The tapes stop where they run out of
  slots --- several here reach six --- which is the saturation of eq.~\eqref{eq:depth}.}
  \label{fig:S5a}
\end{figure}

\begin{figure}[p]
  \centering
  \includegraphics[width=\textwidth]{figS5b_shared_prefix}
  \caption{\textbf{The later two lineages split, the more slots their cells share.} One comparison per
  (internal node, tape): one cell drawn from each side of the split, over 200 trees. Points are the
  mean within a time bin, at the mean split time of that bin; bars are standard errors across trees;
  the line is eq.~\eqref{eq:shared}. All 20 bins sit within $1.04$ standard errors of it, with mixed
  signs.
  \emph{What it says:} the shared prefix is the clock the tree is read from, and it is coarse ---
  cells that split at harvest still share only $4.8$ of 6 slots on average.
  \emph{Two things the shape shows.} The curve \emph{bends} but does not reach the ceiling: by
  eq.~\eqref{eq:depthslope} its slope falls from $5.5$ to $2.9$ slots per unit time, since $47\%$ of
  tapes are full by harvest. So a late split leaves about half as many distinguishing edits per unit
  time as an early one --- fewer, not none. (A prediction that the curve would saturate past $t\approx0.7$
  was wrong, and how much of the late-history limit of Fig.~S4 is saturation rather than simply the
  time remaining before harvest is not yet apportioned.)}
  \label{fig:S5b}
\end{figure}

\begin{figure}[p]
  \centering
  \includegraphics[width=0.86\textwidth]{figS5c_signal_in_slot_order}
  \caption{\textbf{A signal pulse is written into slot order.} Constant editing rate; signal share
  $0.02$ outside the shaded window and $0.30$ inside it ($0.4\le t<0.7$); tapes that filled all six
  slots, $47\%$ of them.
  \textbf{Top:} when each slot was written --- the density $K_{j,6}$ of eq.~\eqref{eq:slotshare}, in
  calendar time. Mean write times run $0.12,0.24,0.36,0.48,0.61,0.73$, and the windows are wide and
  heavily overlapping.
  \textbf{Bottom:} the share of tapes whose slot $j$ holds the signal symbol: simulated
  $2.8,6.3,12.2,16.8,16.2,11.8\%$ against eq.~\eqref{eq:slotshare}'s
  $2.8,6.1,12.0,16.7,16.4,11.9\%$, all within $1.21$ standard errors.
  \emph{How the rows connect:} each dot's height is set by how much of that slot's curve lies inside
  the shaded band. Slot~4's curve lies about half inside, so its share sits about half way between
  $2\%$ and $30\%$. Only $3,15,36,52,51,35\%$ of writes to slots 1--6 fell inside the window, which is
  why a $30\%$ pulse returns as at most $\sim17\%$: \textbf{each slot averages the signal over its own
  write window, so the recorder blurs time rather than losing it.} The share does not fall back to
  zero at slot~6 for three reasons --- the off-share is $0.02$, not $0$; the axis is slot, not time,
  and $35\%$ of slot-6 writes still happened during the pulse; and a full tape stops recording, so the
  late off-period is barely written on these tapes at all.}
  \label{fig:S5c}
\end{figure}

\begin{figure}[p]
  \centering
  \includegraphics[width=0.86\textwidth]{figS5d_programmed_rate}
  \caption{\textbf{The editing rate can be programmed over time.} The \textsc{tv} schedule of
  \S\ref{sec:editval}: $6.47$, $0$, $8.63$, $4.31$ edits per tape per experiment on
  $0$--$0.25$, $0.25$--$0.35$ (the pause, shaded), $0.35$--$0.6$ (the burst) and $0.6$--$1$. No signal
  channel.
  \textbf{Top:} the realised rate, edits falling in a bin $0.025$ wide divided by the tape-time spent
  \emph{open} in that bin, summed over every lineage alive. It sits on the programmed step:
  \emph{exactly zero} edits during the pause, as \S\ref{sec:piecewise} requires, and over the 36 bins
  with editing the deviations have mean $0.17$ and spread $1.00$ standard errors, worst $2.83$ (a
  maximum that large among 36 bins arises by chance about $16\%$ of the time). Open circles divide the
  same edits by \emph{all} tape-time, open or full: they sag --- $8.48\to7.17$ across the burst,
  $3.48\to2.31$ over the last phase --- although the rate is constant within each phase, because tapes
  are filling. That is eq.~\eqref{eq:depthslope}'s two factors pulled apart, and the same distinction
  that panel (b) turns on.
  \textbf{Bottom:} filled slots per tape along lineages, against eq.~\eqref{eq:depth}, within $1.1$
  standard errors at all 41 time points: flat at $1.62$ slots through the pause, steeper through the
  burst, bending as tapes fill, ending at $4.80$ --- the same as the faint constant-rate curve, since
  both schedules spend the same total $\LamT$.
  \emph{What this does not show.} Both rows use the simulator's true edit times, which real data do
  not have. The panel shows the rate can be \emph{imposed}, not that it can be \emph{recovered} from
  sequenced tapes. Recovery runs through the branching times and meets the confound named in
  \texttt{sciphy\_notes.md} \S S4.7: a low early rate and coalescences sitting later --- which higher
  turnover or higher capture both produce --- look alike.}
  \label{fig:S5d}
\end{figure}

\subsection{Why these four}

Taken together the panels cover what the layer has to get right before anything is built on top of it.
Panel~(a) shows the mechanism, and (b) turns it into the quantity every lineage statistic depends on:
how much two cells share, and how coarsely. Panel~(c) is the one that matters for signalling history
--- it shows that a pulse in time survives as a pattern across slots, which is the property the
inference will exploit, and it quantifies the blurring, which is the price. Panel~(d) shows the clock
of \S\ref{sec:lambdatime} working in both directions, including a zero-rate window that needs no
special handling.

Two limits are visible in the figures rather than argued: the shared prefix is only $4.8$ of 6 slots
even for cells that split at harvest (b), and a $30\%$ pulse returns as at most $\sim17\%$ in any slot
(c). Both are recorder limits, not inference limits; what they cost a lineage-resolved signalling
history is the subject of the switch-detectability figure, and what they cost through dropout is the
next layer.
